\begingroup
\renewcommand{\arraystretch}{1.45}
\Needspace{40\baselineskip}
\InterpolationArticle{24}
Durch die seit Artikel 20 dargestellte Untersuchung sind wir dahin gelangt, dass wir bei gegebenen $\mu$
Werten der Funktion $X$ des veränderlichen Winkels $x$, die eine vollständige Periode bilden, eine Funktion $(X')=$
\[
\begin{aligned}
&\gamma+\gamma'\cos x+\gamma''\cos2x+\mathinner{.\,.\,.}+\gamma^m\cos mx\\[5pt]
&\quad+\delta'\sin x+\delta''\sin2x+\mathinner{.\,.\,.}+\delta^m\sin mx
\end{aligned}
\]
angeben können, die allen diesen Werten genügt, wobei die Zahl $m$, welche die Ordnung dieser Funktion ausdrückt, entweder $=\frac12\mu-\frac12$ oder $=\frac12\mu$ ist, je nachdem, ob $\mu$ ungerade oder
gerade ist. Aus den Artikeln 15 und 19 folgt, dass jede andere ähnliche Funktion, die denselben Werten genügt, in der Form $X'+Y\sin\frac12\mu(x-a)$ enthalten sein muss,
wobei $Y$ eine unbestimmte Funktion des Winkels $x$ von der Form $G$ oder $F$ (Artikel 13) bezeichnet, je nachdem, ob $\mu$
ungerade oder gerade ist. Im ersten Fall gibt es außer $X'$ keine andere ähnliche Funktion, die denselben
Werten genügt, außer einer von höherer Ordnung; daher wird $X'$ mit $X$ identisch sein, wenn
feststeht, dass $X$ die $m^{\mathrm{te}}$ Ordnung nicht überschreitet. Im zweiten Fall genügen zwar unendlich viele
andere ähnliche Funktionen der $m^{\mathrm{ten}}$ Ordnung denselben Werten; alle stimmen jedoch darin
überein, dass sie in der Form $X'+h\sin m(x-a)$ enthalten sind, wobei
$h$ eine feste Größe bezeichnet. Wenn also feststeht, dass $X$ die $m^{\mathrm{te}}$ Ordnung nicht überschreitet,
werden sich $X$ und $X'$ jedenfalls höchstens in den Koeffizienten von $\cos mx$ und $\sin mx$
unterscheiden; in der Funktion $X'$ wird jedoch zwischen diesen Koeffizienten die Gleichung
$\gamma^m\sin ma-\delta^m\cos ma=0$ gelten. Ferner ist klar: Ist von anderswoher eine Funktion $X''$ bekannt, die eine ähnliche Form wie $X'$ hat und denselben Werten genügt, so genügt ihnen jede Funktion $X''+Y\sin\frac12\mu(x-a)$; dabei lässt sich $Y$ so bestimmen, dass diese Funktion die $m^{\mathrm{te}}$ Ordnung nicht überschreitet. \Needspace{35\baselineskip}
Zu diesem Zweck bemerken wir: Wird ein Teil von $Y$ durch $K\cos(\kappa x+k)$ dargestellt, wobei $K$
einen festen Koeffizienten, $k$ einen festen Winkel und $\kappa$ eine Zahl bezeichnet, die entweder ganzzahlig ist oder den Bruchteil
$\frac12$ hat (je nachdem, ob $\mu$ gerade oder ungerade ist)---auf diese Form lassen sich die einzelnen Teile bringen, wenn Kosinus und Sinus desselben Winkels zusammengefasst werden---, so wird der entsprechende Teil des Produkts $Y\sin\frac12\mu(x-a)$ durch
\SourceLeaf{311}
\[
\begin{aligned}
&\frac12K\sin((\kappa+\frac12\mu)x+k-\frac12\mu a)\\[5pt]&\qquad{}-\frac12K\sin((\kappa-\frac12\mu)x+k+\frac12\mu a)
\end{aligned}
\]
oder, was dasselbe ist, durch
\[
\begin{aligned}
&L\sin(\lambda x+l)-L\sin((\lambda-\mu)x+l+\mu a)
\end{aligned}
\]
dargestellt, wobei $\lambda$ eine ganze Zahl bezeichnet. Daher darf in $X''$ jeder Teil $L\sin(\lambda x+l)$ durch $L\sin((\lambda-\mu)x+l+\mu a)$ ersetzt werden und dieser auf ähnliche Weise wiederum durch
$L\sin((\lambda-2\mu)x+l+2\mu a)$, dann durch $L\sin((\lambda-3\mu)x+l+3\mu a)$ usw. Um einen entsprechenden Schluss auf Kosinus auszudehnen, genügt die Bemerkung, dass anstelle der Sinus Kosinus auftreten,
wenn nur $l$ durch $l+90^{0}$ ersetzt wird. \Needspace{18\baselineskip}
Daraus folgt: Kommt in $X''$ ein Glied
$L\cos\lambda x$ vor, wobei $\lambda$ eine ganze Zahl größer als $\frac12\mu$ bezeichnet, die sich daher in die Form
$\lambda=\nu\mu\pm\lambda'$ bringen lässt, sodass $\nu$ ganzzahlig und $\lambda'$ nicht größer als $\frac12\mu$ ist, so darf dieses Glied durch
\[
\begin{aligned}
&L\cos(\pm\lambda'x+\nu\mu a)=L\cos\nu\mu a\cos\lambda'x\mp L\sin\nu\mu a\sin\lambda'x
\end{aligned}
\]
ersetzt werden; ebenso darf ein solches Glied $L\sin\lambda x$ durch
\[
\begin{aligned}
&L\sin(\pm\lambda'x+\nu\mu a)=L\sin\nu\mu a\cos\lambda'x\pm L\cos\nu\mu a\sin\lambda'x
\end{aligned}
\]
ersetzt werden. Auf diese Weise wird die Funktion $X''$ offenbar auf eine andere reduziert, die gewiss höchstens von der
$\frac12\mu^{\mathrm{ten}}$ Ordnung ist.

\Needspace{27\baselineskip}
Durch diese Operation geht also $X''$ stets in $X'$ selbst über, wenn $\mu$
ungerade ist; im anderen Fall, in dem $\mu$ gerade ist, dagegen gewiss in eine Funktion, die
sich nur in den Koeffizienten von $\cos mx$ und $\sin mx$ von der Funktion $X'$ unterscheiden
kann. Um aus jenen Koeffizienten die entsprechenden Koeffizienten in $X'$ herzuleiten, kann die allgemeine Methode von Artikel 15 angewandt werden. Aus ihr folgt ohne Schwierigkeit: Sind
die letzten Glieder in jener Funktion $K\cos mx+L\sin mx$, so müssen sie, damit die Funktion $X'$ entsteht, durch
\[
\begin{aligned}
&\quad(K\cos ma+L\sin ma)\cos ma\cos mx\\[5pt]
&+(K\cos ma+L\sin ma)\sin ma\sin mx
\end{aligned}
\]
ersetzt werden. Diese Glieder werden $=K\cos mx+L\sin mx$, wenn zwischen $K$ und $L$ die Gleichung $K\sin ma=L\cos ma$
gilt, wie sich durch Rechnung leicht bestätigt; dann ist jene Funktion also bereits vollständig mit $X'$ identisch.
\endgroup
